\section{Discussion}
\label{sec:discussion}

\subsection{Why Does Error Injection Fail on Embodied Planning?}
\label{sec:structural_hypothesis}

We hypothesize that the transfer failure stems from structural differences between tool-use and embodied planning tasks. Table~\ref{tab:structural} contrasts the two domains across key dimensions.

\begin{table}[t]
\centering
\small
\begin{tabular}{p{2.2cm}p{2.3cm}p{2.3cm}}
\toprule
\textbf{Dimension} & \textbf{Tool-Use} \cite{paladin2025} & \textbf{Embodied} \\
\midrule
Error signal & Discrete (API error codes) & Diffuse (wrong object, wrong room) \\
Observability & Explicit (stack trace) & Silent (no error msg) \\
Recovery scope & Local (retry call) & Global (replan) \\
Feedback & Immediate & Delayed \\
Error locality & Single step & Cascading \\
\bottomrule
\end{tabular}
\caption{Structural comparison: tool-use vs.\ embodied planning. Error injection benefits from discrete, observable, locally recoverable errors.}
\label{tab:structural}
\end{table}

In tool-use, errors produce discrete signals (API error codes, type mismatches) that map directly to specific recovery actions. The agent learns a conditional policy: ``if error type X, then recovery action Y.'' In embodied planning, errors propagate silently through the action sequence---picking the wrong object produces no error message, and recovery requires replanning from an unknown state. The taxonomy labels that help in tool-use (mapping error types to recovery strategies) become noise in embodied planning because there is no clean error$\to$recovery mapping to learn.

\subsection{Alternative Explanation: Recovery-Pairing Confound}
\label{sec:confound}

An important methodological difference between PALADIN and our setup: PALADIN pairs each error trajectory with an explicit recovery demonstration at inference time via retrieval augmentation, contributing an estimated $\sim$30pp of its +57pp gain. Our TaxEI injects standalone errors without paired recoveries. The negative result may therefore stem from missing recovery pairing rather than (or in addition to) task structure. Future work should test recovery-paired injection on ALFWorld to disentangle these factors.

\subsection{Limitations}
\label{sec:limitations}

\paragraph{Scale.} We evaluate on a single environment (ALFWorld, 47 tasks) with one base model (Qwen3-8B). Generalization to other embodied environments (e.g., WebShop) and model scales remains untested.

\paragraph{Statistical power.} Three seeds provide limited statistical power: all comparisons are NS, and the variance gradient ($3.3 < 7.0 < 8.6$pp SD) is directional but underpowered. We acknowledge the tension in arguing that $n{=}3$ is insufficient while reporting only $n{=}3$ ourselves.

\paragraph{Dose-response.} The dose-response ablation uses a single seed, preventing confidence intervals. The threshold observation (10$\to$25\%) requires multi-seed replication.

\paragraph{Data asymmetry.} TaxEI and Random use $2.1\times$ more training data than Vanilla (9{,}029 vs.\ 4{,}358). The fixed-size dose-response ablation (4{,}358 total) partially controls for this, showing that the degradation is not merely a volume effect.

\paragraph{Per-operator analysis.} With $N{=}11$--$12$ tasks per operator, the local/global split is exploratory. A single task flip changes the delta by 8--9pp.

\paragraph{Version discrepancy.} Our evaluation pipeline evolved from v3 (74.5\% vanilla) to v4 (65.3\% vanilla) due to format alignment changes and seed standardization. All reported results use v4.
